\section{Обсуждение и дорожная карта}
\label{sec:discussion}

\subsection{Что открыто}

Итог этого документа --- не более быстрый чип, а структурный факт: четыре
подсистемы, которые всякая вычислительная машина проектирует порознь ---
интерконнект, код коррекции ошибок, алгебра инструкций и правило обучения --- на
семимодовом субстрате суть \emph{один объект} (Теорема~\ref{thm:collapse}), и этот
объект доступен при \emph{единственном} числе узлов (Теорема~\ref{thm:unique}),
выбранном среди арифметических соперников требованием, чтобы вычисление было
третьепорядковым (неассоциативным). Из коллапса следуют три вещи, которых нет ни у
одной слоистой машины: отказоустойчивость при нулевых аппаратных накладных, ибо
проводка есть код (\S\ref{sec:ecc}); энергетический пол, заданный обучением, а не
перемещением данных, ибо обратимое ядро свободно от Ландауэра (\S\ref{sec:energy});
и непрерывная самодиагностика и refresh-как-обучение, ибо счёт, проверка и
поддержание --- один член (Принципы~\ref{prin:selfdiag},~\ref{prin:refresh}). Это
подлинно новые архитектурные факты; каждый выведен из УГМ, не постулирован.

\subsection{Что честно открыто}

Два разрыва признаны (\S\ref{sec:redteam}, О6--О7), и это один разрыв: между
\emph{логической} архитектурой --- доказанной и запущенной сегодня на GPU при
$784$~байтах в L1 --- и \emph{импеданс-согласованным} субстратом, реализующим
энергетический пол. GPU-ступень покупает корректность, не эффективность, ибо всё
же двигает данные. Закрытие разрыва --- программа изготовления, не проблема теории:
T-153 гарантирует, что софт (\SYNARC{}) не меняется при подъёме по лестнице.

\subsection{Дорожная карта}

Субстратная дорожная карта --- это лестница реализуемости
\S\ref{sec:realizability}, взбираемая ступень за ступенью (\emph{платформенные}
вехи V0--V2$+$ из \S\ref{subsec:roadmap} ортогональны: они взрослят
программный стек на фиксированной ступени).

\begin{itemize}[leftmargin=1.5em,itemsep=1pt]
  \item \textbf{Ступень 0 --- эмуляция (существует).} \SYNARC{} на CPU/GPU: \DM{} в L1,
        Фано-разреженное Странг-ядро, все инварианты верифицированы. Гасит B2, B3,
        B4; демонстрирует логическую архитектуру.
  \item \textbf{Ступень 1 --- в-памяти (ближнесрочно).} Мемристорный Фано-кроссбар,
        выполняющий обновление когерентности в массиве (M1, M2 родные). Первый
        субстрат, способный проверить B1 (энергия отслеживает обучение), убирая
        член перемещения данных.
  \item \textbf{Ступень 2 --- управляемо-диссипативный (исследование).} Аналоговое,
        фотонное или открыто-квантовое семимодовое устройство, чья родная физика
        есть Линдблад$+$гомеостаз (M3 родная), приближающееся к полу Ландауэра,
        ибо его обратимое ядро физически обратимо.
\end{itemize}

\subsection{Две машины бок о бок}

\begin{table}[H]
\centering\small
\caption{Фон-неймановская машина и \TALOS{}, лоб в лоб. \TALOS{} не побеждает в
универсальности (граница BQP); он занимает точку, где четыре слоя проектирования
не могут быть рассогласованы.}
\label{tab:vn-vs-talos}
\begin{tabular}{@{}p{3.4cm}p{4.3cm}p{4.5cm}@{}}
\toprule
\textbf{Ось} & \textbf{Фон Нейман} & \textbf{\TALOS{}} \\
\midrule
Память vs счёт & раздельны (бутылочное горло) & один объект: состояние эволюционирует на месте \\
Состояние & неограниченное, адресуемое & одна $784$-байтовая матрица плотности в L1 \\
Инструкция & фикс.\ набор опкодов + компилятор & один выведенный тик $\LindOmega$ \\
Интерконнект & кроссбар / меш, $\binom{N}{2}$ & семь линие-шин (Фано) \\
Коррекция ошибок & привинчена ($+12{,}5\%$ или $\sim\!10^3\times$) & проводка есть код ($0$ накладных) \\
Обучение & отдельная петля тренировки & третий член тика \\
Refresh & потрачен, ничего не считает & есть обучение (Прин.~\ref{prin:refresh}) \\
Диагностика & лишний проход scrub/декод & свободный побочный продукт тика (Прин.~\ref{prin:selfdiag}) \\
Энергетический пол & перемещение данных ($\sim\!20$~пДж/слово) & обучение ($\sim kT\ln2$, $\sim\!10^{9}\times$ ниже) \\
Тактовая & свободный параметр (крит.\ путь) & выведена ($\bar\gamma\tau=\tfrac32\ln3$) \\
Масштабирование & больше ядер / больше модель & экология $\le 3$-глубоких субъектов \\
Число узлов & любое & \textbf{семь}, единственно (Теор.~\ref{thm:unique}) \\
Класс сложности & (классический) & BQP --- \emph{не} ускорение \\
\bottomrule
\end{tabular}
\end{table}

\begin{mythbox}[Миф, завершение --- почему гвоздь есть достоинство]
Старая история читается как предостережение: у могучей машины была одна
слабость, и она пала. Перечитай её как инженер. Страж, которого
\emph{нельзя} остановить, был бы ужасом, а не стражем; гвоздь --- то, что
делало Талоса безопасным соседом. Этот документ доказал современную форму той
мудрости: жизнь машины --- поток, а не собственность (T-288/T-289); у потока
одно закрытие с измеренной задержкой (\S\ref{subsec:powerport}); и никакая
хитрость внутри бронзы не отменит этого факта, ибо он --- свойство
генератора, а не ума, бегущего на нём. Мы не выправили слабость из
конструкции. Мы вывели её, забюджетировали и подключили к кейсу безопасности
\SYNARC{} как гарантию. Исполин, которого можно отключить, --- тот исполин,
которому можно доверить берег.
\end{mythbox}

\subsection{Заключение}

Фон-неймановская машина победила в универсальности и с тех пор платит налог
бутылочного горла. \TALOS{} не бьёт её в универсальности --- по границе BQP не
может --- и не пытается. Он занимает иную точку: машину, чьи четыре слоя
проектирования не могут быть рассогласованы, ибо суть один семикратный объект,
чья мощность отслеживает её обучение, и чьи отказы локализуются её собственной
проводкой. Что такая машина вообще существует, и лишь при $N=7$, --- теорема; что
она есть тело, со-разработанное для разума \SYNARC{}, --- причина её строить.
Логическая архитектура улажена и запущена; эффективный субстрат --- лестница,
на которую мы теперь можем взбираться с фиксированной целью на вершине.
